\chapter{Radio front-end impairments}\label{ch:impairments}

Real transmitters and receivers distort the signal before any propagation effect. This chapter
models the impairments implemented in \file{phy/channel/impairments.py}. \Cref{fig:impair-const}
shows their signatures on a 16-QAM constellation.

\begin{figure}[htbp]
\centering
\begin{tikzpicture}
\begin{groupplot}[group style={group size=4 by 1,horizontal sep=4mm,y descriptions at=edge left},
  constaxis,width=4.1cm,height=4.1cm]
\nextgroupplot[title={AWGN only}]
  \addplot[scatterpts=accent] table[x=i,y=q]{\datafile{const_awgn}};
\nextgroupplot[title={phase noise}]
  \addplot[scatterpts=accent3] table[x=i,y=q]{\datafile{const_pn}};
\nextgroupplot[title={IQ imbalance}]
  \addplot[scatterpts=accent4] table[x=i,y=q]{\datafile{const_iq}};
\nextgroupplot[title={Saleh PA, 4\,dB IBO}]
  \addplot[scatterpts=accent2] table[x=i,y=q]{\datafile{const_pa}};
\end{groupplot}
\end{tikzpicture}
\caption{16-QAM at $E_s/N_0=26$\,dB through \blk{phase_noise} (150\,Hz linewidth at
1\,MBd), \blk{iq_imbalance} (1.5\,dB, $8^\circ$) and \blk{power_amplifier} (Saleh). Phase
noise smears points along circles; IQ imbalance turns the square grid into a parallelogram;
the PA compresses and rotates the outer points.}
\label{fig:impair-const}
\end{figure}

% ------------------------------------------------------------------
\section{Carrier frequency and phase offsets}
A local oscillator at $f_c+\Delta f$ with phase $\phi_0$ turns the received envelope into
\begin{equation}
  y[n]=x[n]\,e^{\,j(2\pi\Delta f\,n/f_s+\phi_0)} .
\end{equation}
\blk{carrier_offset} applies this rotation (\prm{cfo_hz}, \prm{phase_deg}, or a uniformly random
phase per frame). Synchronization (\cref{ch:sync}) estimates and removes it.

% ------------------------------------------------------------------
\section{Oscillator phase noise}\label{sec:phase-noise}
A free-running oscillator accumulates phase errors. The standard model, and the one used for
lasers in \cref{ch:optical}, is a Wiener process: the instantaneous frequency deviation is
white, so the phase performs a random walk. In discrete time,
\begin{equation}
  \phi[n]=\phi[n-1]+\varepsilon[n],\qquad \varepsilon[n]\sim\N\!\Big(0,\ \frac{2\pi\Delta\nu}{f_s}\Big).
  \label{eq:wiener}
\end{equation}
\begin{proposition}
The field $e^{j\phi(t)}$ of a Wiener phase with increment variance $2\pi\Delta\nu|\tau|$ has a
Lorentzian spectrum of full width at half maximum $\Delta\nu$.
\end{proposition}
\begin{proof}
$\phi(t+\tau)-\phi(t)\sim\N(0,2\pi\Delta\nu|\tau|)$, and $\E e^{jX}=e^{-\Var X/2}$ for a Gaussian
$X$, so the autocorrelation is $R(\tau)=e^{-\pi\Delta\nu|\tau|}$. Its Fourier transform is
\[
  S(f)=\frac{\Delta\nu/(2\pi)}{f^2+(\Delta\nu/2)^2},
\]
which falls to half its peak at $f=\pm\Delta\nu/2$.
\end{proof}
\codeexcerpt{phase_noise}{Wiener phase noise (\blk{phase_noise}).}
The phase variance grows linearly with time. What matters for a receiver is therefore the
product $\Delta\nu\,T$ of linewidth and symbol duration: carrier recovery must track the phase
faster than it wanders.

% ------------------------------------------------------------------
\section{IQ imbalance}
In a quadrature down-converter the Q branch may have a gain $g\ne1$ and a phase error $\varphi$
relative to the I branch. Writing the two branch outputs and recombining them gives the
widely-linear model
\begin{equation}
  y=K_1x+K_2x\conj+d,\qquad K_1=\frac{1+g\,e^{-j\varphi}}{2},\qquad K_2=\frac{1-g\,e^{j\varphi}}{2},
  \label{eq:iq}
\end{equation}
where $d$ is a DC offset. The conjugate term creates a mirror image of the signal. Its
strength is measured by the image rejection ratio $\mathrm{IRR}=|K_1|^2/|K_2|^2$; for $g=1.5$\,dB
and $\varphi=8^\circ$ it is about $19$\,dB. On a constellation it shears the square grid
(\cref{fig:impair-const}).
\codeexcerpt{iq_imbalance}{IQ imbalance and DC offset (\blk{iq_imbalance}).}

% ------------------------------------------------------------------
\section{Power amplifier nonlinearity}
A memoryless band-pass nonlinearity acts on the envelope $x=re^{j\theta}$ as
\begin{equation}
  y=A(r)\,e^{\,j(\theta+\Phi(r))},
\end{equation}
where $A(r)$ is the AM/AM curve and $\Phi(r)$ the AM/PM curve. The input is first scaled so that
its mean power lies \prm{input_backoff_db} below the saturation level $A_{\mathrm{sat}}=1$
(input back-off, IBO). Three models are implemented:
\begin{align}
  \text{Rapp (solid state):}\quad & A(r)=\frac{r}{\big(1+r^{2p}\big)^{1/(2p)}},\qquad\Phi=0,\\
  \text{Saleh (TWTA):}\quad & A(r)=\frac{r}{1+\beta_a r^2},\qquad \Phi(r)=\frac{\alpha_\phi r^2}{1+\beta_\phi r^2},\\
  \text{ideal clipper:}\quad & A(r)=\min(r,1),\qquad\Phi=0.
\end{align}
The Saleh AM/AM curve is divided by $\alpha_a$ so that all models have unit small-signal gain;
the default constants are those of Saleh's original TWT fit~\cite{saleh1981}. The smoothness
$p$ of the Rapp model interpolates between a soft compression ($p=1$) and the ideal clipper
($p\to\infty$).

\begin{figure}[htbp]
\centering
\begin{tikzpicture}
\begin{groupplot}[group style={group size=2 by 1,horizontal sep=18mm},width=6.8cm,height=5.2cm,
  xlabel={input amplitude $r$},xmin=0,xmax=2]
\nextgroupplot[ylabel={output amplitude $A(r)$},ymin=0,ymax=1.2,legend pos=south east,title={AM/AM}]
  \addplot[th=black!40,dashed,domain=0:1.2] {x};
  \addplot[th=accent] table[x=r,y=rapp1]{\datafile{pa_curves}};
  \addplot[th=accent3] table[x=r,y=rapp2]{\datafile{pa_curves}};
  \addplot[th=accent4] table[x=r,y=rapp5]{\datafile{pa_curves}};
  \addplot[th=black!70] table[x=r,y=clip]{\datafile{pa_curves}};
  \addplot[th=accent2] table[x=r,y=saleh]{\datafile{pa_curves}};
  \legend{linear,Rapp $p{=}1$,Rapp $p{=}2$,Rapp $p{=}5$,clipper,Saleh}
\nextgroupplot[ylabel={phase shift $\Phi(r)$ [deg]},ymin=0,ymax=25,title={AM/PM (Saleh)}]
  \addplot[th=accent2] table[x=r,y=saleh_pm_deg]{\datafile{pa_curves}};
\end{groupplot}
\end{tikzpicture}
\caption{PA characteristics measured by driving the \blk{power_amplifier} block with an amplitude
ramp (back-off chosen so that the block's input scaling is unity).}
\label{fig:pa}
\end{figure}
\codeexcerpt{power_amplifier}{Back-off scaling and the AM/AM, AM/PM models.}

\begin{pitfall}
The nonlinearity generates spectral regrowth (out-of-band emission) and an in-band error that
is \emph{not} Gaussian. For OFDM, whose envelope is nearly Gaussian (\cref{sec:papr}), an IBO
of several dB is needed. Compare the \blk{power_meter} PAPR with the back-off.
\end{pitfall}

% ------------------------------------------------------------------
\section{Analog-to-digital conversion}
\blk{quantizer} models a uniform mid-rise ADC with $b$ bits on each of I and Q. The full scale
$A$ is either given or set to \prm{loading_factor}$\times$ the RMS value $\sigma_x$ of each
component. The step is $\Delta=2A/2^b$. For a signal that does not clip, the quantization error is
approximately uniform with variance $\Delta^2/12$, giving
\begin{equation}
  \mathrm{SQNR}=\frac{\sigma_x^2}{\Delta^2/12}=\frac{3\cdot2^{2b}}{(A/\sigma_x)^2}
  \ \Rightarrow\ \mathrm{SQNR}_{\mathrm{dB}}\approx6.02\,b+4.77-20\log_{10}\frac{A}{\sigma_x}.
\end{equation}
With the default loading factor of 4 (clipping probability about $6\times10^{-5}$ for a Gaussian
component), an 8-bit ADC gives about 41\,dB.
\codeexcerpt{quantizer}{Uniform mid-rise quantization with clipping.}

% ------------------------------------------------------------------
\section{Timing offset}
\blk{timing_offset} delays the signal by $d=\lfloor d\rfloor+\mu$ samples. The integer part is a
shift. The fractional part uses a Blackman-windowed sinc interpolator of 33 taps,
$h[n]=\sinc(n-16-\mu)\,w[n]/\sum_n\sinc(n-16-\mu)\,w[n]$, which is accurate as long as the signal
is oversampled. The frame length is preserved, so a delay pushes the tail out of the frame; the
\blk{preamble} block appends guard symbols for this reason.
